\documentclass[11pt]{article}

\usepackage[margin=1in]{geometry}
\usepackage{amsmath, amssymb}
\usepackage{parskip}
\usepackage[hidelinks]{hyperref}
\pagestyle{plain}

\newenvironment{reviewerquote}
  {\begin{quote}\itshape}
  {\end{quote}}

\title{Response to reviewers\\
  \large Introspection Dynamics with Mutation in Additive Games}
\author{Harry Foster, Vince Knight, Sebastian Krapohl}
\date{\today}

\begin{document}

\maketitle

Dear Professor Zaccour,

We thank both reviewers for their positive assessments and for a further
set of careful comments, all of which we have addressed. We respond to
each comment below; reviewer text is set in italics, and our responses
follow in upright type.

\section*{Reviewer 1}

\begin{reviewerquote}
Page 3: ``At each step, one player \(i\) is selected uniformly at random,
and draws an alternative action \(b\) uniformly from the \(L_i-1\) actions
other than \(a_i\).'' This sentence can be confusing because it sounds as
if a player cannot mutate to the current action. And, to fully describe
the process, you should also define the transition probability for when
\(b=a_i\).
\end{reviewerquote}

Done, on both counts. The process is now described as an explicit
two-stage update: with probability \(\mu_{ib}\) the selected player
mutates to action \(b\), and the target may be their current action, in
which case the state does not change; otherwise, with probability
\(1-M_i\), they introspect, drawing an alternative uniformly from the
other \(L_i-1\) actions and switching with probability given by the Fermi
function. The holding probability \(T^{\mathrm{intro}}_{\mathbf{a},
\mathbf{a}}\) is now displayed as a numbered equation immediately after
the transition probability, with a sentence stating which events it
collects (mutation to the current action, and introspection without
switching).

\begin{reviewerquote}
Page 4: the \(m_i\) notation for the stationary distribution could be
understood as a scalar, when it's a vector. Consider using boldface.
\end{reviewerquote}

Done. The stationary distribution is now written \(\mathbf{m}_i\)
throughout, consistent with the boldface state \(\mathbf{a}\).

\begin{reviewerquote}
Page 4: in the last equation (after equation 6),
\(T_{\mathbf{a},\mathbf{a}}\) doesn't seem obvious to me (at least, I had
to verify it on the side). Or perhaps I missed something.
\end{reviewerquote}

The reviewer did not miss anything; the expression was stated without
justification. It now follows by inspection: the holding probability is
the numbered equation added in response to the first comment, and the
proof of Theorem 1 states that the expression follows from it together
with the diagonal entries
\(K_i(a_i,a_i)=1-\sum_{b\neq a_i}K_i(a_i,b)\).

\begin{reviewerquote}
Page 5: in Theorem 2's proof, you refer to equation 5 when I think you
meant the equation next to 5 (which is not numbered).
\end{reviewerquote}

Quite right: the intended reference was the formula for the off-diagonal
entries of the selection kernel \(K_i\), which does not carry the
\(1/N\) factor. That formula is now a numbered equation and the proof
cites it.

\begin{reviewerquote}
Page 5: ``By Theorem 1 the stationary marginal of player \(i\)'s action
is \(m_i\), so the long-run probability that player \(i\) cooperates is
\(m_i(C) = p_i\).'' I think Theorem 1 doesn't even need to be invoked
here, right?
\end{reviewerquote}

Right again. The proof now gives the direct argument: additivity
decouples player \(i\) from the rest of the population, so their action
evolves as a two-state Markov chain in its own right, with kernel
\((1-\tfrac{1}{N})I+\tfrac{1}{N}K_i\), which has the same stationary
distribution as \(K_i\). Theorem 1 is no longer invoked.

\begin{reviewerquote}
Page 5: ``Additivity makes a player's choice forgetful:\ldots'' maybe
make it clear here again that is both additivity AND having only 2
actions that makes a player's choice `forgetful'.
\end{reviewerquote}

Done, and the comment prompted us to correct a related error. The sentence
now reads ``Additivity and the restriction to two actions together make a
player's choice forgetful''. In addition, the property is now named and
stated where it is introduced, rather than only being pointed at ahead of
Theorem 2: it is the property that the two rows of \(K_i\) coincide, so
that the probability that a selected player's next action is \(C\) does
not depend on the action they currently hold. We had justified its failure
for \(L_i\ge3\) by saying that the probability of retaining the current
action depends on that action, which is mistaken, since this is equally
true of the two-action case, where the retention probabilities are
\(p_i\) and \(1-p_i\). The text now gives the correct reason in both
directions: with a single alternative the two switching probabilities are
complementary, by \(\phi_i(x)+\phi_i(-x)=1\), and it is this that aligns
the rows; for \(L_i\ge3\) the rows always differ, since identical rows
would force \(g_i\) to be constant, and equating a diagonal entry with the
off-diagonal entries in its column would then require \(L_i=2\). The
closed form of Theorem 2 is thus special to two actions in a strict
sense, and not merely generically.

\begin{reviewerquote}
Page 16: I assume you use some standard correlation measure (Pearson for
binary variables?) but it would be nice to specify which one is exactly
used.
\end{reviewerquote}

Done. Both the body text and the caption of the figure now state that
\(\mathrm{corr}(a_i,a_j)\) is the Pearson correlation coefficient of the
binary actions under the exact stationary distribution (for indicator
variables, the phi coefficient), which is what the accompanying source
code computes.

\section*{Reviewer 2}

\begin{reviewerquote}
First about Equation (15). I believe Eq.\ (15) corresponds to the
sentence ``[\ldots] each unit contributed changes the marginal value of
the next by a factor \(w > 0\) [\ldots]'', only when the object
\(X(\mathbf{a})\) is an integer (i.e., the contributions \(\alpha_i\) are
integers). For continuous values of contributions, there has to be some
other normalization that is not \((w-1)\). Therefore, if you want to
stick to continuous values of alpha, I suggest that you mention that this
formula is just a canonical extension of the `synergy' model. Or, perhaps
you would like to derive the other formula from first principle or
mention that alpha is restricted to integers only.
\end{reviewerquote}

The reviewer is correct, and we have adopted their first suggestion. The
text now states that the geometric-sum reading requires the total
contribution \(X(\mathbf{a})\) to be an integer, as it is in both of our
examples (\(\alpha_i=1\) for the top row of the figure and
\(\alpha_i=i\) for the correlation panels), and that for non-integer
contributions we regard the payoff formula as the canonical analytic
extension of the synergy and discounting model, with the normalisation
\(w-1\) retained so that the linear game is still recovered as
\(w\to1\).

\begin{reviewerquote}
Finally, I was wondering what the purpose of studying this model is,
broadly speaking (in evolutionary terms). For example, in standard games
(let's just say in additive games or the non-linear public goods game
that the authors study) would players ever learn to be more exploratory
(i.e., would this mutational model itself evolve and spread in a
population that does only pure introspection selection)? Does
incorporating mutations in decisions have a fitness advantage? Can this
be the case or is this model simply dominated as a strategy if the
co-player is not exploring? I would have liked if the authors commented
on that aspect. It is also quite easy to perform analysis on this since
the authors already have payoff formulas.
\end{reviewerquote}

We thank the reviewer for this question, and we agree that the closed
forms make the analysis tractable. Indeed, the additive structure gives a
sharp preliminary answer. By the additive decomposition and the product
measure, a player's expected stationary payoff is affine in their own
cooperation probability \(p_i\), with slope \(-\delta_i\), whilst the
co-players' marginals are entirely beyond their influence. In a social
dilemma (\(\delta_i>0\)) symmetric exploration raises \(p_i\) above the
pure-introspection value \(\phi_i(\delta_i)<\tfrac{1}{2}\), so it
strictly lowers the explorer's own payoff whatever the co-players do:
unbiased exploration is dominated, confirming the reviewer's conjecture,
and in a form stronger than conjectured, since by the decoupling it does
not matter whether the co-players explore. Two further observations
follow from
the same formulas. First, what is dominated is exploration towards the
disfavoured action rather than noise as such: mutation biased towards
defection pushes \(p_i\) below \(\phi_i(\delta_i)\) and outperforms pure
introspection. Second, exploration towards cooperation benefits every
co-player, so it is altruistic, and its evolution inherits the structure
of the underlying dilemma. Beyond additivity the decoupling fails, and
our Figure 4 shows that cooperation can peak at a strictly positive
mutation rate, so there exploration may plausibly carry an individual
advantage.

We believe a proper treatment of this question, with heritable update
rules competing on the slower evolutionary timescale, and the
non-additive case examined, is a study in its own right, and, mindful
that the manuscript already grew substantially in the previous round, we
have not added it to the paper. Instead the conclusion now closes with a
new paragraph that does two things. It first makes the evolutionary
interpretation of our results explicit: the mutation rates model
trembles, experimentation, and behavioural noise, and on this reading
the results describe how noisy self-reflection shapes long-run
behaviour, keeping every action in play, confining each player strictly
inside the pure conventions however strong selection becomes, and
leaving cooperation at a mutation-selection balance. It then poses the
reviewer's question, whether such noise would itself be selected for if
update rules were heritable and fitness were the expected stationary
payoff, citing the exploration dynamics of Traulsen et al.\ (2009), and
notes that the closed forms make it tractable and that it is a natural
direction for future work.

\vspace{2em}

Thank you again for the constructive reports.

Harry Foster, Vince Knight, and Sebastian Krapohl

\end{document}
